\documentclass[11pt,a4paper]{article}
\usepackage[utf8]{inputenc}
\usepackage[T1]{fontenc}
\usepackage[english]{babel}
\usepackage{amsmath,amssymb,amsthm}
\usepackage{graphicx}
\usepackage{hyperref}
\usepackage{listings}
\usepackage{xcolor}
\usepackage{geometry}
\usepackage{float}
\usepackage{caption}
\usepackage{subcaption}
\usepackage{algorithm}
\usepackage{algpseudocode}
\usepackage{booktabs}
\usepackage{siunitx}

\geometry{margin=2.5cm}

\title{Multi-Point Tactile Sensor Force Regression and Classification Models}
\author{MagTec Models Research Team}
\date{\today}

\hypersetup{
    colorlinks=true,
    linkcolor=blue,
    filecolor=magenta,      
    urlcolor=cyan,
    pdftitle={Multi-Point Tactile Sensor Models},
}

\lstset{
    language=Python,
    basicstyle=\ttfamily\small,
    keywordstyle=\color{blue},
    stringstyle=\color{red},
    commentstyle=\color{green!60!black},
    numbers=left,
    numberstyle=\tiny\color{gray},
    stepnumber=1,
    numbersep=5pt,
    backgroundcolor=\color{gray!10},
    frame=single,
    breaklines=true,
    showstringspaces=false
}

\begin{document}

\maketitle

\begin{abstract}
This document describes the development and implementation of machine learning models for multi-point tactile sensor systems. The models enable force regression (Fx, Fy, Fz components) and classification tasks (location and stretch level) using magnetic field sensor data. The system utilizes a combination of normal force and shear force datasets, achieving high accuracy across all tasks: Fz regression RMSE of \SI{0.0854}{\newton} (R² = 0.9898), Fx regression RMSE of \SI{0.2021}{\newton} (R² = 0.9827), Fy regression RMSE of \SI{0.1512}{\newton} (R² = 0.9900), location classification accuracy of 99.59\%, and stretch classification accuracy of 99.30\%.
\end{abstract}

\tableofcontents
\newpage

\section{Introduction}

Multi-point tactile sensing systems require accurate force estimation and spatial localization capabilities. This document presents a comprehensive machine learning framework for processing magnetic field sensor data from a 15-taxel StretchMagTec sensor array to predict contact forces and classify contact locations and material stretch levels.

\subsection{System Overview}

The tactile sensing system consists of:
\begin{itemize}
    \item \textbf{Sensor Array}: 15 magnetic field sensors (taxels) arranged in a 3×5 grid
    \item \textbf{Force/Torque Sensor}: FT-300S sensor providing ground truth force measurements
    \item \textbf{Data Collection}: Robotic manipulation system (Franka Emika Panda) for controlled data acquisition
    \item \textbf{ML Models}: Random Forest regressors and classifiers for force prediction and classification
\end{itemize}

\section{Mathematical Framework}

\subsection{Feature Extraction}

The magnetic field sensor provides a 3D vector measurement for each taxel, resulting in a feature vector $\mathbf{m} \in \mathbb{R}^{15 \times 3} = \mathbb{R}^{45}$ for each time sample.

\subsubsection{Raw Features}

The raw feature vector $\mathbf{x}_{\text{raw}}$ is constructed by flattening the magnetic field tensor:
\begin{equation}
\mathbf{x}_{\text{raw}} = \text{vec}(\mathbf{M}) = [m_{1,x}, m_{1,y}, m_{1,z}, \ldots, m_{15,x}, m_{15,y}, m_{15,z}]^T \in \mathbb{R}^{45}
\end{equation}
where $\mathbf{M} \in \mathbb{R}^{15 \times 3}$ is the magnetic field matrix and $m_{i,j}$ represents the $j$-th component (x, y, z) of the $i$-th taxel.

\subsubsection{Magnitude Features}

An alternative feature representation uses the magnitude (L2 norm) of each taxel's magnetic field vector:
\begin{equation}
\mathbf{x}_{\text{mag}} = [\|\mathbf{m}_1\|_2, \|\mathbf{m}_2\|_2, \ldots, \|\mathbf{m}_{15}\|_2]^T \in \mathbb{R}^{15}
\end{equation}
where $\|\mathbf{m}_i\|_2 = \sqrt{m_{i,x}^2 + m_{i,y}^2 + m_{i,z}^2}$ is the Euclidean norm of the $i$-th taxel's magnetic field vector.

\subsubsection{Normalized Magnitude Features}

For improved robustness, magnitude features can be normalized per sensor:
\begin{equation}
\mathbf{x}_{\text{mag,norm}} = \left[\frac{\|\mathbf{m}_1\|_2 - m_{1,\min}}{m_{1,\max} - m_{1,\min}}, \ldots, \frac{\|\mathbf{m}_{15}\|_2 - m_{15,\min}}{m_{15,\max} - m_{15,\min}}\right]^T
\end{equation}
where $m_{i,\min}$ and $m_{i,\max}$ are the minimum and maximum magnitudes observed for taxel $i$ across the training dataset.

\subsection{Data Preprocessing}

\subsubsection{Baseline Removal}

For shear force regression, baseline removal is applied to isolate force variations from initial sensor offsets. The baseline $\bar{\mathbf{f}}_0$ is computed from the first 10\% of samples in each sequence:
\begin{equation}
\bar{f}_{0,k} = \frac{1}{N_0} \sum_{n=1}^{N_0} f_{k,n}, \quad k \in \{x, y, z\}
\end{equation}
where $N_0 = \lceil 0.1 \cdot N \rceil$ and $N$ is the sequence length. The baseline-removed force is:
\begin{equation}
\tilde{f}_{k,n} = f_{k,n} - \bar{f}_{0,k}
\end{equation}

\textbf{Important}: For Fz regression with normal forces, baseline removal is \textit{not} applied, as absolute force values are required.

\subsubsection{Displacement Filtering}

High-frequency noise is filtered using displacement-based filtering. The displacement $\Delta_n$ between consecutive samples is computed:
\begin{equation}
\Delta_n = \|\mathbf{x}_n - \mathbf{x}_{n-1}\|_2
\end{equation}
Samples exceeding the 95th percentile threshold are removed:
\begin{equation}
\mathcal{M} = \{n : \Delta_n < \text{percentile}(\{\Delta_i\}, 95)\}
\end{equation}

\subsubsection{Force Clipping}

For regression targets, forces are clipped to valid ranges:
\begin{align}
\tilde{f}_x &\leftarrow \text{clip}(\tilde{f}_x, -2.0, 2.0) \quad \text{[N]} \\
\tilde{f}_y &\leftarrow \text{clip}(\tilde{f}_y, -2.0, 2.0) \quad \text{[N]} \\
f_z &\leftarrow \min(f_z, 3.0) \quad \text{[N]}
\end{align}

\subsection{Force Regression Models}

\subsubsection{Fz Regression}

The Fz regressor $\mathcal{R}_{Fz}: \mathbb{R}^{45} \rightarrow \mathbb{R}$ predicts normal force magnitude:
\begin{equation}
\hat{f}_z = \mathcal{R}_{Fz}(\mathbf{x})
\end{equation}

\textbf{Dataset}: Normal forces only (test42)
\begin{itemize}
    \item 30 sequences per location × 11 locations = 330 sequences per stretch level
    \item Total: 990 sequences across 3 stretch levels
    \item \textbf{No baseline removal} (absolute force values)
\end{itemize}

\textbf{Model}: Random Forest Regressor with 200 estimators

\textbf{Performance}:
\begin{align}
\text{RMSE} &= \SI{0.0854}{\newton} \\
R^2 &= 0.9898
\end{align}

\subsubsection{Fx Regression}

The Fx regressor $\mathcal{R}_{Fx}: \mathbb{R}^{45} \rightarrow \mathbb{R}$ predicts shear force in the x-direction:
\begin{equation}
\hat{f}_x = \mathcal{R}_{Fx}(\mathbf{x})
\end{equation}

\textbf{Dataset}: Shear forces only (test51), x+ and x- sequences
\begin{itemize}
    \item 20 sequences per location per direction × 10 locations × 2 directions = 400 sequences per stretch
    \item Total: ~1200 sequences across 3 stretch levels
    \item \textbf{Baseline removal applied}
    \item \textbf{Fx inversion}: Due to action-reaction principle, Fx is inverted for x-direction movements
\end{itemize}

\textbf{Model}: Random Forest Regressor with 200 estimators

\textbf{Performance}:
\begin{align}
\text{RMSE} &= \SI{0.2021}{\newton} \\
R^2 &= 0.9827
\end{align}

\subsubsection{Fy Regression}

The Fy regressor $\mathcal{R}_{Fy}: \mathbb{R}^{45} \rightarrow \mathbb{R}$ predicts shear force in the y-direction:
\begin{equation}
\hat{f}_y = \mathcal{R}_{Fy}(\mathbf{x})
\end{equation}

\textbf{Dataset}: Shear forces only (test51), y+ and y- sequences
\begin{itemize}
    \item 20 sequences per location per direction × 10 locations × 2 directions = 400 sequences per stretch
    \item Total: ~1200 sequences across 3 stretch levels
    \item \textbf{Baseline removal applied}
\end{itemize}

\textbf{Model}: Random Forest Regressor with 200 estimators

\textbf{Performance}:
\begin{align}
\text{RMSE} &= \SI{0.1512}{\newton} \\
R^2 &= 0.9900
\end{align}

\subsection{Classification Models}

\subsubsection{Location Classification}

The location classifier $\mathcal{C}_{\text{loc}}: \mathbb{R}^{45} \rightarrow \{0, 1, \ldots, 10\}$ maps magnetic field features to contact locations:
\begin{equation}
\hat{\ell} = \mathcal{C}_{\text{loc}}(\mathbf{x})
\end{equation}
where $\ell \in \{0, \ldots, 9\}$ represents contact points 1-10, and $\ell = 10$ represents no contact.

\textbf{Dataset}: Combined normal + shear forces
\begin{itemize}
    \item Normal: 330 sequences per stretch (30 per location)
    \item Shear: ~800 sequences per stretch (20 per location × 4 directions)
    \item Combined: ~1130 sequences per stretch
    \item Total: ~3390 sequences across 3 stretch levels
\end{itemize}

\textbf{Feature Method}: Raw features (45 dimensions) or Magnitude features (15 dimensions)

\textbf{Model}: Random Forest Classifier with 200 estimators, class\_weight='balanced'

\textbf{Performance}:
\begin{align}
\text{Accuracy} &= 99.59\% \\
\text{Features} &= 45 \text{ (raw)} \text{ or } 15 \text{ (magnitude)}
\end{align}

\subsubsection{Stretch Classification}

The stretch classifier $\mathcal{C}_{\text{stretch}}: \mathbb{R}^{45} \rightarrow \{0, 1, 2\}$ predicts material stretch level:
\begin{equation}
\hat{s} = \mathcal{C}_{\text{stretch}}(\mathbf{x})
\end{equation}
where $s \in \{0, 1, 2\}$ corresponds to \{0\%, 10\%, 20\%\} stretch levels.

\textbf{Dataset}: Combined normal + shear forces (same as location classification)

\textbf{Feature Method}: Raw features (45 dimensions)

\textbf{Model}: Random Forest Classifier with 200 estimators

\textbf{Performance}:
\begin{equation}
\text{Accuracy} = 99.30\%
\end{equation}

\section{Model Architecture}

\subsection{Random Forest Algorithm}

All models utilize Random Forest (RF) algorithms, which combine multiple decision trees through ensemble learning.

\subsubsection{Regression}

For regression tasks, the RF regressor aggregates predictions from $B$ trees:
\begin{equation}
\hat{y} = \frac{1}{B} \sum_{b=1}^{B} T_b(\mathbf{x})
\end{equation}
where $T_b: \mathbb{R}^{d} \rightarrow \mathbb{R}$ is the $b$-th decision tree.

Each tree $T_b$ is trained on a bootstrap sample of the training data and selects features randomly at each split to reduce overfitting.

\subsubsection{Classification}

For classification tasks, the RF classifier uses majority voting:
\begin{equation}
\hat{y} = \arg\max_{c} \sum_{b=1}^{B} \mathbf{1}[T_b(\mathbf{x}) = c]
\end{equation}
where $\mathbf{1}[\cdot]$ is the indicator function and $c$ is a class label.

\subsection{Hyperparameters}

\begin{table}[H]
\centering
\begin{tabular}{@{}ll@{}}
\toprule
\textbf{Parameter} & \textbf{Value} \\
\midrule
Number of estimators & 200 \\
Max depth & None (unlimited) \\
Random state & 42 \\
Class weight (classification) & 'balanced' \\
\bottomrule
\end{tabular}
\caption{Random Forest hyperparameters}
\label{tab:hyperparams}
\end{table}

\section{Training Procedure}

\subsection{Data Splitting}

Sequences are split at the sequence level (not sample level) to prevent data leakage:
\begin{equation}
\mathcal{D}_{\text{train}} = 70\% \text{ of sequences}, \quad \mathcal{D}_{\text{test}} = 30\% \text{ of sequences}
\end{equation}

The random seed is fixed (seed=42) for reproducibility.

\subsection{Feature Normalization}

Features are standardized using StandardScaler:
\begin{equation}
\mathbf{x}_{\text{norm}} = \frac{\mathbf{x} - \boldsymbol{\mu}}{\boldsymbol{\sigma}}
\end{equation}
where $\boldsymbol{\mu}$ and $\boldsymbol{\sigma}$ are the mean and standard deviation computed on the training set.

\subsection{Outlier Removal}

Outliers are removed using Median Absolute Deviation (MAD) method:
\begin{equation}
\text{MAD} = \text{median}(|X_i - \text{median}(X)|)
\end{equation}
Sequences with modified z-scores exceeding a threshold ($z_{\text{thresh}} = 3.0$) are removed:
\begin{equation}
z_i = \frac{0.6745 \cdot (X_i - \text{median}(X))}{\text{MAD}}
\end{equation}

\section{Performance Metrics}

\subsection{Regression Metrics}

\subsubsection{Root Mean Squared Error (RMSE)}
\begin{equation}
\text{RMSE} = \sqrt{\frac{1}{N} \sum_{i=1}^{N} (y_i - \hat{y}_i)^2}
\end{equation}

\subsubsection{Coefficient of Determination (R²)}
\begin{equation}
R^2 = 1 - \frac{\sum_{i=1}^{N} (y_i - \hat{y}_i)^2}{\sum_{i=1}^{N} (y_i - \bar{y})^2}
\end{equation}
where $\bar{y} = \frac{1}{N}\sum_{i=1}^{N} y_i$ is the mean of the true values.

\subsection{Classification Metrics}

\subsubsection{Accuracy}
\begin{equation}
\text{Accuracy} = \frac{\text{TP} + \text{TN}}{\text{TP} + \text{TN} + \text{FP} + \text{FN}}
\end{equation}
where TP, TN, FP, FN are true positives, true negatives, false positives, and false negatives, respectively.

\section{Code Architecture and Usage}

\subsection{Repository Structure}

The codebase is organized as follows:

\begin{lstlisting}[language=bash]
Mag_Tec_Models/
├── src/
│   ├── training/
│   │   ├── train_best_models.py          # Main training script
│   │   ├── train_separate_forces.py      # Separate force regressors
│   │   ├── train_location_all_methods.py # Location classification
│   │   └── train_combined_shear_normal.py # Utility functions
│   ├── franka_controller/
│   │   ├── franka_skin_test_multiple_points.py  # Normal force collection
│   │   └── franka_skin_test_shear_forces.py     # Shear force collection
│   └── ...
├── data/
│   └── Multiple_Points/
│       ├── 2.5mm_single_test42/         # Normal forces dataset
│       └── shear_forces_test51/          # Shear forces dataset
│           └── best_models/
│               └── best_models/
│                   └── models/           # Trained models
├── docs/
│   └── model_documentation.tex           # This document
└── README.md
\end{lstlisting}

\subsection{Command Line Interface}

\subsubsection{Training Best Models}

To train all models with optimal configurations:

\begin{lstlisting}[language=bash]
python3 src/training/train_best_models.py \
    --normal-dir data/Multiple_Points/2.5mm_single_test42 \
    --shear-dir data/Multiple_Points/shear_forces_test51 \
    --run-label best_models \
    --remove-outliers
\end{lstlisting}

\subsubsection{Training Separate Force Regressors}

To train Fz, Fx, and Fy regressors separately:

\begin{lstlisting}[language=bash]
python3 src/training/train_separate_forces.py \
    --normal-dir data/Multiple_Points/2.5mm_single_test42 \
    --shear-dir data/Multiple_Points/shear_forces_test51 \
    --run-label separate_forces \
    --remove-outliers
\end{lstlisting}

\subsubsection{Training Location Classification}

To train location classifiers with all feature methods:

\begin{lstlisting}[language=bash]
python3 src/training/train_location_all_methods.py \
    --normal-dir data/Multiple_Points/2.5mm_single_test42 \
    --shear-dir data/Multiple_Points/shear_forces_test51 \
    --run-label location_all_methods \
    --remove-outliers
\end{lstlisting}

\subsection{Data Collection}

\subsubsection{Normal Forces Collection}

\begin{lstlisting}[language=bash]
python3 src/franka_controller/franka_skin_test_multiple_points.py
\end{lstlisting}

This script collects normal force data at 10 contact points plus no-touch sequences across 3 stretch levels.

\subsubsection{Shear Forces Collection}

\begin{lstlisting}[language=bash]
python3 src/franka_controller/franka_skin_test_shear_forces.py
\end{lstlisting}

This script collects shear force data with movements in 4 directions (x+, x-, y+, y-) at each contact point.

\subsection{Database Structure}

\subsubsection{HDF5 File Format}

Data is stored in HDF5 format with the following structure:

\begin{lstlisting}[language=Python]
h5_file/
├── presses/
│   ├── press_000/
│   │   ├── forces          # [N, 6] array (Fx, Fy, Fz, Tx, Ty, Tz)
│   │   ├── stretchmagtec   # [N, 15, 3] array (magnetic field data)
│   │   ├── timestamps      # [N] array
│   │   └── attributes:
│   │       ├── label       # String identifier
│   │       ├── offset      # Location ID ('1'-'10' or 'no_touch')
│   │       └── stretch_label # '000pct', '010pct', '020pct'
│   ├── press_001/
│   └── ...
└── labels/                 # [M] array of sequence labels
\end{lstlisting}

\subsubsection{Data Organization}

\begin{itemize}
    \item \textbf{Normal Forces}: One HDF5 file per stretch level
    \begin{itemize}
        \item \texttt{*\_stretch\_000pct.h5}
        \item \texttt{*\_stretch\_010pct.h5}
        \item \texttt{*\_stretch\_020pct.h5}
    \end{itemize}
    
    \item \textbf{Shear Forces}: One HDF5 file per stretch level
    \begin{itemize}
        \item \texttt{*\_shear\_stretch\_000pct.h5}
        \item \texttt{*\_shear\_stretch\_010pct.h5}
        \item \texttt{*\_shear\_stretch\_020pct.h5}
    \end{itemize}
\end{itemize}

\subsubsection{Sequence Metadata}

Each sequence contains:
\begin{itemize}
    \item \textbf{Location}: Contact point ID ('1' through '10') or 'no\_touch'
    \item \textbf{Stretch Level}: Material stretch percentage (0\%, 10\%, 20\%)
    \item \textbf{Direction}: For shear forces, movement direction ('x+', 'x-', 'y+', 'y-')
    \item \textbf{Sampling Rate}: 100 Hz
    \item \textbf{Duration}: ~6 seconds per sequence
\end{itemize}

\section{Model Files}

\subsection{Saved Models}

Trained models are saved in Joblib format:

\begin{itemize}
    \item \texttt{fz\_regressor.joblib} - Fz force regressor
    \item \texttt{fx\_regressor.joblib} - Fx force regressor
    \item \texttt{fy\_regressor.joblib} - Fy force regressor
    \item \texttt{location\_classifier.joblib} - Location classifier (raw features)
    \item \texttt{stretch\_classifier.joblib} - Stretch level classifier
\end{itemize}

\subsection{Scalers}

Feature scalers are saved separately:

\begin{itemize}
    \item \texttt{scaler\_fz.joblib} - Scaler for Fz regression features
    \item \texttt{scaler\_fx.joblib} - Scaler for Fx regression features
    \item \texttt{scaler\_fy.joblib} - Scaler for Fy regression features
    \item \texttt{location\_scaler.joblib} - Scaler for location classification features
    \item \texttt{stretch\_scaler.joblib} - Scaler for stretch classification features
\end{itemize}

\section{Results Summary}

\subsection{Force Regression Performance}

\begin{table}[H]
\centering
\begin{tabular}{@{}lccccc@{}}
\toprule
\textbf{Component} & \textbf{Dataset} & \textbf{RMSE [N]} & \textbf{R²} & \textbf{Train Seq} & \textbf{Test Seq} \\
\midrule
Fz & Normal only & 0.0854 & 0.9898 & 693 & 297 \\
Fx & Shear only (x±) & 0.2021 & 0.9827 & 832 & 357 \\
Fy & Shear only (y±) & 0.1512 & 0.9900 & 831 & 357 \\
\bottomrule
\end{tabular}
\caption{Force regression performance}
\label{tab:force_regression}
\end{table}

\subsection{Classification Performance}

\begin{table}[H]
\centering
\begin{tabular}{@{}lccccc@{}}
\toprule
\textbf{Task} & \textbf{Features} & \textbf{Accuracy [\%]} & \textbf{Train Seq} & \textbf{Test Seq} & \textbf{Classes} \\
\midrule
Location & Raw (45) & 99.59 & 893 & 383 & 11 \\
Location & Magnitude (15) & 99.59 & 893 & 383 & 11 \\
Stretch & Raw (45) & 99.30 & 893 & 383 & 3 \\
\bottomrule
\end{tabular}
\caption{Classification performance}
\label{tab:classification}
\end{table}

\section{Conclusion}

The developed models achieve high accuracy across all tasks:
\begin{itemize}
    \item Fz regression: RMSE of \SI{0.0854}{\newton} with R² = 0.9898
    \item Fx regression: RMSE of \SI{0.2021}{\newton} with R² = 0.9827
    \item Fy regression: RMSE of \SI{0.1512}{\newton} with R² = 0.9900
    \item Location classification: 99.59\% accuracy
    \item Stretch classification: 99.30\% accuracy
\end{itemize}

The system demonstrates robust performance suitable for real-time tactile sensing applications in robotic manipulation.

\section{Acknowledgments}

This work was supported by the MagTec Models research project. The authors acknowledge the use of the Franka Emika Panda robotic platform and StretchMagTec sensor hardware.

\end{document}

